\documentclass[11pt,a4paper]{article}
\usepackage[T1]{fontenc}
\usepackage{lmodern,amsmath,amssymb,amsthm,mathtools,microtype,booktabs,array,xcolor}
\usepackage[margin=25mm,top=23mm,bottom=25mm,headheight=15pt]{geometry}
\usepackage{fancyhdr,enumitem,listings}
\lstset{language=Python,basicstyle=\ttfamily\footnotesize,columns=fullflexible,keepspaces=true,showstringspaces=false,breaklines=true}
\usepackage[colorlinks=true,linkcolor=blue!35!black,citecolor=blue!40!black,urlcolor=blue!42!black]{hyperref}
\usepackage{bookmark}
\definecolor{ink}{HTML}{153C48}
\definecolor{teal}{HTML}{23635D}
\newtheorem{theorem}{Theorem}[section]
\newtheorem{lemma}[theorem]{Lemma}
\newtheorem{proposition}[theorem]{Proposition}
\newtheorem{corollary}[theorem]{Corollary}
\theoremstyle{remark}\newtheorem{remark}[theorem]{Remark}
\numberwithin{equation}{section}
\setlength{\emergencystretch}{1.5em}
\pagestyle{fancy}\fancyhf{}\lhead{\footnotesize Spectral sampling, independent sets, and Lucas sequences}
\rhead{\footnotesize Research manuscript}\cfoot{\small\thepage}
\renewcommand{\headrulewidth}{.3pt}
\hypersetup{pdftitle={Sampling at Bessel Zeros, Fixed-Order Independence Sequences, and Lucas Divisibility}}
\begin{document}
\begin{center}
{\Large\bfseries\color{ink} Sampling at Bessel Zeros, Fixed-Order Independence\\[3pt] Sequences, and Lucas Divisibility\par}
\vspace{7pt}
{\small Research manuscript prepared for Keio\enspace $\cdot$\enspace 20 November 2026}
\end{center}
\begin{abstract}
We prove the coefficient summability problem of Berdysheva and Berens in every dimension $d\ge3$, with a bound uniform over normalized admissible profiles. Integration of a positive Fourier representation supplies the spectral moment of order $d-1$ needed to sample at spherical Bessel zeros. In the associated discrete Tur\'an problem we prove that the first zero uniquely maximizes in dimensions 2 through 28 and the third zero uniquely maximizes in dimensions 29 through 32. We also determine the limiting maximizing radius and index in high dimension. Two independent applications concern finite structures. A bound on joint extensions of independent sets forbids a prescribed valley in the fixed-order Roller-Coaster conjecture for infinitely many parameter pairs. For very well-covered graphs it implies unimodality through independence number 35 and log-concavity through 33. Connected chordal bipartite examples admit arbitrarily many separated peaks. In Sun's Conjecture~4.4 we classify every allowed parameter. Zero is the sole exception to the two universal divisibility assertions. For every nonzero admissible parameter and every starting index, explicit progressions supply infinitely many witnesses for both Lucas recurrences.
\end{abstract}

\section{The problems and their proof thresholds}
The main analytic question concerns a passage from a positive continuous spectrum to infinitely many discrete coefficients. The corresponding coefficient sum had been established in dimension two. The original question asks for all profiles in every higher dimension~\cite[Open Problem~1]{BB16}. We prove a uniform bound for this entire class. The same spherical Bessel nodes occur in two other questions in the source. For those extremal questions we give exact low-dimensional answers and an asymptotic formula for large dimensions.

The graph question fixes both the order and the independence number before prescribing a permutation of coefficient ranks~\cite[Conj.~1.5]{LMroller}. A single forbidden ordering at fixed parameters refutes the conjecture. We give a restriction on three ranks, a hierarchy for several ranks, and their consequences. For the arithmetic question~\cite[Conj.~4.4]{Sun2019}, the requirement is an infinite set of indices for each of two recurrences. A witness for one congruence need not witness the other.

Three patterns in the supplied corpus motivate the proofs. Entry~185 derives an extremum from a differential equation and a signed Wronskian. We compare reflected oscillatory solutions and bound the integrated area remaining after each zero. Entry~216 extracts small local witnesses from sets in a graph. Here the local data count vertices available separately and pairs available together. The prime-power lifts in entries~199 and~204 motivate a third question about prescribing a scalar matrix modulo its own exponent. The corpus entries concern different conclusions. In particular, entry~204 already contains the five-sequence Chen--Guo--Wang application, and it is excluded from the results here~\cite{VMRC}. Sections~2 and 3 address the analytic problems, Sections~4 and 5 concern graphs, and Section~6 treats the divisibility question.

\section{A uniform summability theorem for the Berdysheva--Berens coefficients}
\label{sec:sampling}
Let $d\ge3$. In the notation of Berdysheva and Berens, $\mathcal K_\pi^{(d)}$ consists of continuous real $\varphi$ supported on $[0,\pi]$, with $\varphi(0)=1$, for which the following even transform is nonnegative:
\[
 \psi(s)=2\int_0^\pi t^{d-1}\varphi(t)h_d(st)\,dt,\qquad
 h_1(z)=\cos z,\quad h_{d+1}(z)=h_d(z)+d\int_0^1 v^{d-1}h_d(vz)\,dv.
\]
This is the source's normalization in its equation~(1). Equivalently,
$f(x)=\varphi(\|x\|_1)$ is positive definite on $\mathbb R^d$~\cite[Sec.~3]{BB16}. In particular $|\varphi|\le1$. Put $l=d-2$, let $r_n$ be the positive zeros of $j_l$ in increasing order, and let $s_n=r_n/\pi$. With $j_l(x)=\sqrt{\pi/(2x)}J_{l+1/2}(x)$, the defining integral for the biorthogonal coefficients~\cite[pp.~69--70]{BB16} is
\begin{equation}
 A_{n;d}=\frac{1}{C_{n;d}}\int_0^\pi t^{d-1}\varphi(t)h_d(s_nt)\,dt,\qquad
 C_{n;d}=s_n^{3-d}\frac{\pi^3}{2} j_{l+1}(r_n)^2 .
 \label{eq:An}
\end{equation}
In particular $A_{n;d}=\psi(s_n)/(2C_{n;d})$ with the factor 2 in our definition of $\psi$. The source also writes its sampling formula with a different overall factor. A fixed positive normalization has no effect on summability.

\begin{theorem}[Berdysheva--Berens Open Problem 1]\label{thm:sampling}
For each $d\ge3$ there exists $K_d<\infty$ such that
\[
   \sup_{\varphi\in\mathcal K_\pi^{(d)}}\,
   \sum_{n\ge1} A_{n;d}(\varphi)\le K_d .
\]
\end{theorem}

\subsection{The positive spectral moment}
The auxiliary question is which weighted mass of $\psi$ is available before applying any inversion formula at zero. Use the Fourier convention $\widehat f(x)=(2\pi)^{-d}\int f(y)e^{-i x\cdot y}\,dy$.
\begin{lemma}\label{lem:mass}
For every admissible $\varphi$,
\[
  \int_0^\infty s^{d-1}\psi(s)\,ds
       =2^{d-1}\Gamma(d/2)^2.
\]
\end{lemma}
\begin{proof}
Let $\Delta$ be the simplex $t_i\ge0$, $\sum_{i=1}^d t_i=1$, with measure $dt_1\cdots dt_{d-1}$, of total mass $1/(d-1)!$. The positive B-spline representation of the Fourier transform~\cite[Theorem A]{BB16} reads
\begin{equation}
 \widehat f(x)=(2\pi)^{-d}\int_{\Delta}
      \psi\!\left(\sqrt{\smash[b]{\textstyle\sum_i t_i x_i^2}}\right)\,dt .
 \label{eq:simplex}
\end{equation}
For distinct knots $x_i^2$, this follows from the source's B-spline formula and the Hermite--Genocchi identity~\cite{deboor}. Specifically, take a $(d-1)$-fold antiderivative of $u\mapsto\psi(\sqrt u)$ on the interval of knots. Its divided difference equals both the integral of that derivative against the source's spline and its integral over $\Delta$ at the barycentric combination of the knots. The exceptional set of repeated knots has Lebesgue measure zero in $\mathbb R^d$.

Bochner's theorem gives a finite representing measure of total mass $f(0)=1$. Since $f$ is integrable, this measure has continuous density $\widehat f$ in our convention. Thus $\int\widehat f=1$. In \eqref{eq:simplex} the integrand is nonnegative, so Tonelli applies without any advance claim about its integral. For an interior $t$ put $y_i=\sqrt{t_i}x_i$ and use polar coordinates. We obtain
\[
 1=\frac{|\mathbb S^{d-1}|}{(2\pi)^d}
       \left(\int_\Delta\prod_{i=1}^d t_i^{-1/2}\,dt\right)
       \int_0^\infty s^{d-1}\psi(s)\,ds .
\]
The simplex integral is $\Gamma(1/2)^d/\Gamma(d/2)$ and $|\mathbb S^{d-1}|=2\pi^{d/2}/\Gamma(d/2)$. Substitution gives the assertion. This also recovers the formal normalization at zero in the inverse transform and proves that its integral exists.
\end{proof}

\subsection{Sampling at the weights of the biorthogonal system}
A continuous bounded function $P$ whose distributional Fourier transform is supported in $[-a,a]$ obeys the following weighted upper sampling inequality at every $\delta$-separated set $S\subset\mathbb R$. For $b\ge0$,
\begin{equation}
  \sum_{s\in S}(1+|s|)^b|P(s)|
  \le C(a,b,\delta)\int_{\mathbb R}(1+|v|)^b|P(v)|\,dv .
  \label{eq:sampling}
\end{equation}
Indeed, take a Schwartz convolution kernel $K$ with Fourier multiplier identically one on a neighborhood of $[-a,a]$. Then $P=K*P$. The inequality $1+|s|\le(1+|s-v|)(1+|v|)$ transfers the weight to the integrand. In each unit interval there are at most $1+1/\delta$ points of $S-v$, so rapid decay bounds $\sum_{s\in S}(1+|s-v|)^b|K(s-v)|$ uniformly in $v$. This proves \eqref{eq:sampling}.

\begin{proof}[Proof of Theorem~\ref{thm:sampling}]
The recurrence for $h_d$ represents it as a Fourier transform of a positive measure supported on $[-1,1]$. This follows inductively from the two atoms representing the cosine, because dilation by $0\le v\le1$ preserves the support bound. Its mass is $h_d(0)=2^{d-1}$. Integrating the dilated kernels shows that $\psi$ has frequency support in $[-\pi,\pi]$ and $\|\psi\|_\infty\le 2^d\pi^d/d$. This bound on $[-1,1]$ and Lemma~\ref{lem:mass} on its complement give a uniform bound on $\int_{\mathbb R}(1+|s|)^{d-1}\psi(s)\,ds$.

At the zeros, $j_{l+1}(r_n)=-j_l'(r_n)$, and fixed-order Bessel asymptotics give $|r_nj_l'(r_n)|\to1$ as $n\to\infty$~\cite[Secs.~10.17, 10.47, 10.51]{BBDLMF}. All the zeros are simple, so $C_{n;d}^{-1}\le c_d(1+s_n)^{d-1}$ for every $n$. The nodes are separated since successive zeros are asymptotically $\pi$ apart and all finite initial gaps are positive. Apply \eqref{eq:sampling} to $P=\psi$, $S=\{s_n\}$ and $b=d-1$, and substitute \eqref{eq:An}.
\end{proof}

\section{Extremal spherical Bessel zeros in the discrete Tur\'an problem}
\label{sec:bb-ratios}
For $l=d-2$ let $r_n$ be the $n$th positive zero of $j_l(x)=\sqrt{\pi/(2x)}J_{l+1/2}(x)$. Write
\[
 D_l(r)=\int_0^r j_l(x)\,dx,\qquad
 R_l(r_n)=\frac{-j_l'(r_n)}{D_l(r_n)},\qquad
 S_l=\sup_{n\ge1}R_l(r_n).
\]
Open Problem 3 of \cite[p.~71]{BB16} asks for precisely $S_l$ when $d=4,6$
or $d\ge8$. Open Problem 2 asks for the corresponding discrete extremal
constant $\alpha_\pi^{(d)}$ on the source's class $\mathcal A_\pi^{(d)}$ of
normalized profiles with nonnegative spectral coefficients. Its reduction
\cite[eqs.~(7)--(9)]{BB16} is exact:
\begin{equation}
 \alpha_\pi^{(d)}=
 2\pi^d\frac{(d-1)!(d-2)!}{(2d-2)!}\,S_l .
 \label{eq:bb-original}
\end{equation}
Indeed the objective, divided by this prefactor, is an average of $R_l(r_n)$
with nonnegative weights $A_nD_l(r_n)/(d-2)!$ of sum one.
Here $m_{n;d}$ denotes the biorthogonal mode defined on p.~69 of \cite{BB16}, and the original objective is $2\int_0^\pi t^{d-1}\varphi(t)\,dt$. Each single mode, normalized to take value one at the origin, is admissible. Positivity $D_l(r_n)>0$ for every $n$
is the Cooke property recorded in \cite{BB14N}.
Berdysheva's 2014 question asked if the first zero always maximizes the
ratio \cite{BB14N,BB14}. In 2016 the authors reported failure of the
stronger adjacent-peak inequality in high dimensions, and predicted eventual
decrease for each fixed dimension \cite[after Theorem~4]{BB16}.

\subsection{Reflection, energy, and a finite cutoff}
Set $L=l(l+1)$, $w=xj_l$ and $q=1-L/x^2$. The equation is
$w''+qw=0$. The regular solution is positive with positive slope near
zero and, while $w>0$ and $q\le0$, is convex. Its first positive zero
therefore exceeds $\sqrt L$. On $x>\sqrt L$, set $E=w'^2+qw^2$.
Differentiation and the fixed-order sine asymptotics give
\begin{equation}
 E'=q'w^2\ge0,\qquad (E/q)'=-q'w'^2/q^2\le0,\qquad
 E(\infty)=1,\qquad q\le E\le1.
 \label{eq:bb-energy}
\end{equation}

Successive full lobe areas $\int_{r_k}^{r_{k+1}}|j_l|$ strictly decrease.
To see this, reflect the preceding lobe about $r=r_k$ for $k\ge2$.
Choose signs for $u(t)=\pm w(r-t)$ and $v(t)=\mp w(r+t)$ so both begin
with the same positive slope. While both are positive,
$(v'u-u'v)'=(q(r-t)-q(r+t))uv\le0$.
Thus $v/u\le1$, and the outgoing lobe ends at or before the reflected
incoming one. Since $r+t>r-t$, integration with the spherical weight $1/x$
makes the inequality of areas strict, also when $l=0$.
At every odd zero the tail starts with a negative lobe. Its integrals
alternate with decreasing absolute values and tend to zero
by the fixed-order sine asymptotic. The Mellin integral
\cite[eq.~10.22.43]{BBDLMF} consequently gives
\begin{equation}
 D_l(r_{2k+1})>C_l:=\int_0^\infty j_l(t)\,dt
 =\frac{\sqrt\pi\,\Gamma((l+1)/2)}{2\Gamma((l+2)/2)},\qquad
 0<R_l(r_{2k+1})<\frac{1}{C_l r_{2k+1}} .
 \label{eq:bb-tail}
\end{equation}
The second bound uses $|w'(r)|\le1$ at the zero.
Thus every maximizing radius lies below
$B_l=1/(C_lR_l(r_1))$. By Sturm comparison with $q\le1$ the zeros have
spacing at least $\pi$. In particular
\[
 S_l=\max_{0\le k\le \lfloor (B_l-r_1)/(2\pi)\rfloor}R_l(r_{2k+1}).
\]
This explicit finite cutoff holds for every $d$. A bound depending only on $l$ follows by taking $X=2\sqrt{L+1}$. On $[X,X+4\pi]$ one has $q\ge3/4$, so Sturm comparison gives a negative-slope zero $s$ in this interval. At that zero $-w'(s)\ge\sqrt3/2$. The one-integration estimate $|D_l(s)-C_l|\le2/(sq(s))$, proved below, yields $R_l(s)\ge(\sqrt3/2)/[(X+4\pi)C_l+8/3]$. Thus any maximizing radius satisfies
\[
 r_*<\frac{2}{\sqrt3}\left(2\sqrt{L+1}+4\pi+\frac{8}{3C_l}\right).
\]
This bounds the number of candidate zeros by $O(l+1)$.

\subsection{Exact maximizers in dimensions four and six}
For $d=4,6$ the first zero is the unique maximizer, giving the sharp value
\[
 S_{d-2}=\frac{-j'_{d-2}(r_1)}{D_{d-2}(r_1)},\qquad
 \alpha_\pi^{(d)}=2\pi^d\frac{(d-1)!(d-2)!}{(2d-2)!}
 \frac{-j'_{d-2}(r_1)}{D_{d-2}(r_1)} .
\]
The unique maximizing profile in the source's notation is
$2\pi(d-2)!\,m_{1;d}/D_{d-2}(r_1)$, by concentration of the weights in
\eqref{eq:bb-original}.
For the proof, the table lists rational bracket centers; each radius
has error $10^{-5}$. The ratio intervals are multiplied by $10^5$.
\[
\begin{array}{c|rrr|rr}
 l&r_1&r_2&r_3&10^5R_l(r_1)&10^5R_l(r_3)\\ \hline
 2&5.76346&9.09501&12.32294&(17092,17094)&(9258,9260)\\
 4&8.18256&11.70491&15.03966&(15278,15280)&(9854,9856)\\
 30&36.62838&41.64301&46.01804&(7095,7099)&(7231,7235)
\end{array}
\]
By zero spacing, every later positive peak has $r\ge r_1+2\pi$.
For $l=2$ the first bracket gives $r_1>5$ and $C_2=\pi/4$, hence
$1/[C_2(r_1+2\pi)]<4/33<0.17092$. For $l=4$ use $r_1>8$ and
$C_4=3\pi/16$ to get $1/[C_4(r_1+2\pi)]<16/126<0.15278$.
Equation \eqref{eq:bb-tail} excludes all subsequent odd indices;
even-indexed ratios are negative. For $l=30$ the same
certificate exhibits a strict reversal already at the third zero.

Here are reproducible rational instructions for the table. At any listed
center or bracket endpoint $x$ start from $T_0=x^l/(2l+1)!!$ and use
\[
 T_{k+1}=-\frac{x^2 T_k}{2(k+1)(2l+2k+3)} .
\]
For $j_l,j_l',D_l$ multiply $T_k$ respectively by
$1,(l+2k)/x,x/(l+2k+1)$. For $x\le47$ and $k\ge90$ each of the three
absolute terms decreases (the derivative multiplies the ratio by at most
$1+2/(l+180)$). The sums through $k=90,91$ therefore give the two
rational endpoints of an enclosure. On the left and right endpoints of
the three brackets the signs of $j_l$ are respectively $(+,-),(-,+),(+,-)$.
At each center compute the derivative and integral enclosures, then enlarge
them by $10^{-5}/3$ and $10^{-5}$, respectively, before dividing with the
indicated minus sign. These error bounds follow from
$j_l(x)=(2i^l)^{-1}\int_{-1}^1e^{ixt}P_l(t)\,dt$ and $|P_l|\le1$,
which yield $|j_l''|\le1/3$ and $|j_l|\le1$.
To check the indices, no root lies at or before $\sqrt L$. The first
bracket contains a downward crossing, hence an odd-indexed root. Three
roots up to its right endpoint $b$ would span at least $2\pi/\sqrt{q(b)}$
by comparison; each row satisfies $b-\sqrt L<2\pi/\sqrt{q(b)}$.
Successive displayed brackets are closer than $2\pi$ even at opposite
endpoints. An intervening extra zero would require two gaps of at least
$\pi$. This verifies both the enclosures and their labels.

\subsection{The first change of maximizing index}
For every integer $2\le d\le32$, the sharp value in
\eqref{eq:bb-original} and its unique extremizing profile are
\[
 S_{d-2}=R_{d-2}(r_{n_*}),\qquad
 \varphi_*(t)=\frac{2\pi(d-2)!}{D_{d-2}(r_{n_*})}m_{n_*;d}(t),
 \qquad n_*=\begin{cases}1&2\le d\le28,\\3&29\le d\le32.\end{cases}
\]
Thus dimension $29$ is the first failure of the first-zero assertion.
The cutoff inequality \eqref{eq:bb-tail} controls the infinite range of
indices. The rational certificate below supplies all finite comparisons.

Use the three series enclosures defined above with truncations $145,146$,
now for $x\le70$. The ratio bound for consecutive terms beyond this
cutoff is less than $1/10$, including the derivative multiplier.
Define $c_0=157/100$, $c_1=1$, and $c_l=c_{l-2}(l-1)/l$ for $l\ge2$.
Then $0<c_l\le C_l$ by the gamma identity and $\pi>3.14$.
For each $l=0,\ldots,30$ do the following rational operations.

\begin{enumerate}
\item Starting at $t=\max(1,l)$, scan the integer intervals
$[t,t+1],\ldots,[69,70]$. Retain every interval on which the signs of
$j_l$ change from positive to negative. Bracket signs are obtained from
the two consecutive partial sums. The starting sign is positive, and
there is no earlier zero: for $l=0$ use $j_0=\sin(x)/x$, and for $l>0$
use $l<\sqrt L$. No unit interval contains two zeros.
\item Bisect each retained interval $18$ times by the same rational sign
test. If $u$ and $e$ are the final midpoint and half-width, let
$[v,w]$ and $[y,z]$ be the series enclosures for $j_l'(u)$ and $D_l(u)$.
Associate to the bracket the interval
\[
 [L_i,U_i]=
 \left[\frac{-w-e/3}{z+e},\frac{-v+e/3}{y-e}\right].
\]
Here $i=1,\ldots,k$ numbers only the negative-slope zeros in the scan,
so $i$ represents zero $2i-1$. Every test point has a definite rational
sign and $y>e$, $-w>e/3$ in every retained interval.
\item Set $i_*=1$ for $0\le l\le26$ and $i_*=2$ for $27\le l\le30$.
The arithmetic checks are
\[
 L_{i_*}-\max\bigl(\{(70c_l)^{-1}\}\cup
                    \{U_i:1\le i\le k,\ i\ne i_*\}\bigr)>1/12000.
\]
For a short independent check, the following are integer parts of $10^6$
times the smallest left side in each indicated group:
\[
\begin{array}{c|rrrrrrr}
 l&0\!:\!5&6\!:\!10&11\!:\!15&16\!:\!20&21\!:\!25&26&27\!:\!30\\ \hline
 \text{floor}&45756&21405&10414&4479&918&380&101 .
\end{array}
\]
All operations before taking these integer parts are rational.
At every unscanned odd zero $r>70$, equation \eqref{eq:bb-tail} gives
$R_l(r)<(70c_l)^{-1}$, proving uniqueness over the full infinite sequence.
\end{enumerate}

\subsection{Eventual decrease and the limiting maximizer}
At a negative-slope zero $r$, write $A=-w'(r)$. By \eqref{eq:bb-energy},
$\sqrt{q(r)}\le A\le1$. Put $h(x)=1/(xq(x))=x/(x^2-L)$, so that
$h'<0$ and $h''>0$ beyond $\sqrt L$. Twice integrating
$j_l=-h w''$ by parts gives
\[
 \int_r^\infty j_l=-A h(r)-\int_r^\infty w(x)h''(x)\,dx,\qquad
 \left|\int_r^\infty w h''\right|\le\frac{-h'(r)}{\sqrt{q(r)}}.
\]
For fixed $l$ this implies $A=1+O_l(r^{-2})$ and
$D_l(r)=C_l+r^{-1}+O_l(r^{-2})$. Hence
\[
 R_l(r)=\frac{1}{C_l r}-\frac{1}{C_l^2 r^2}+O_l(r^{-3}).
\]
Successive negative-slope zeros are separated by at least $2\pi$,
and by at most $4\pi$ once $q\ge1/2$, by Sturm comparison.
The first term decreases between them by a quantity bounded below by
a positive multiple of $r^{-2}$. The difference in the second terms
and the total remainder are $O_l(r^{-3})$. This proves the eventual
strict decrease predicted after Theorem~4 of \cite{BB16}, in every dimension.

Finally, as $m=l+1/2\to\infty$,
\begin{equation}
 \sqrt m\, S_l\longrightarrow \frac{2}{\sqrt\pi\,3^{3/4}},
 \qquad \frac{r_*}{m}\longrightarrow\sqrt{3/2}
 \quad\hbox{for every choice of a maximizing zero }r_* .
 \label{eq:bb-limit}
\end{equation}
Indeed $C_l\sim\sqrt{\pi/(2m)}$. Write $z=x/m$, $Q(z)=1-z^{-2}$.
The oscillatory large-order expansions for $J_m(mz)$ and its derivative
\cite[Secs.~10.19(ii), 10.20(i)]{BBDLMF} yield, uniformly on compact
subintervals of $1<z<\infty$,
\[
\begin{aligned}
 w(mz)&=Q(z)^{-1/4}\cos\Phi+O(m^{-1}),\\
 w'(mz)&=-Q(z)^{1/4}\sin\Phi+O(m^{-1}),\\
 \Phi&=m\bigl(\sqrt{z^2-1}-\arccos(1/z)\bigr)-\pi/4 .
\end{aligned}
\]
Squaring cancels the phase and gives
$E(mz)=\sqrt{Q(z)}+O(m^{-1})$. Indeed, one integration of $-h w''$ gives $h(r)w'(r)+\int_r^\infty h'w'$ for the tail. Since $|w'|\le1$ and $h$ decreases to zero, $|D_l(r)-C_l|\le2h(r)=2/(rq(r))$ at a zero.
Consequently, uniformly for odd zeros in such a compact $z$-interval,
\[
 \sqrt m\,R_l(r)=\sqrt{2/\pi}\,Q(z)^{1/4}/z+o(1).
\]
Sturm comparison provides bounded gaps between odd zeros in a fixed
compact subinterval of this oscillatory region, so that each interior
$z$ is approximated by $r/m$. For the lower edge $r\le(1+\delta)m$ use
$D_l(r)>C_l$, $r>\sqrt L$, and $E(r)\le E((1+\delta)m)$. The resulting
limsup of $\sqrt m R_l$ is $O(\delta^{1/4})$ as $\delta\downarrow0$.
For $r\ge Bm$, \eqref{eq:bb-tail} bounds the limsup by
$\sqrt{2/\pi}/B$. These two bounds exclude loss of maximizing zeros at
the edges. The continuous limiting function has its unique maximum at
$z=\sqrt{3/2}$, as differentiation shows. This proves both limits in
\eqref{eq:bb-limit}.

The index itself therefore obeys a simple dimension law. Writing $n_*$ for
any maximizing index, one has
\[
 \lim_{l\to\infty}\frac{n_*}{l+1/2}
 =\frac{1/\sqrt 2-\arctan(1/\sqrt 2)}{\pi}
 =0.0291658030\ldots .
\]
To verify this, fix $z>1$ and count zeros in $[\sqrt L,mz]$.
For $0<\delta<z-1$, the part before $m(1+\delta)$ has at most
$1+(m(1+\delta)-\sqrt L)/\pi$ zeros by the spacing bound. On a fixed
partition of $[1+\delta,z]$, apply Sturm comparison on each scaled
subinterval using the minimum and maximum of $q$ there. The zero count
is bounded by the corresponding lengths times $\sqrt{q_{\min}}/\pi$ and
$\sqrt{q_{\max}}/\pi$, up to an additive constant per subinterval.
Divide by $m$, pass to the limit, refine the partition and send
$\delta\downarrow0$. The normalized count tends to
$\pi^{-1}\int_1^z\sqrt{1-u^{-2}}\,du$. Squeezing the maximizing radius
between $m(\sqrt{3/2}\pm\varepsilon)$ in \eqref{eq:bb-limit} and evaluating
the integral proves the assertion.

\section{A fixed-order obstruction to the Roller-Coaster Conjecture}
\label{sec:fixedroller}
Let $s_i$ count the independent $i$-sets of a finite simple graph, with
$s_0=1$. A graph is well-covered if every maximal independent set has the same
cardinality $q=\alpha(G)$. Levit and Mandrescu's overhauled Roller-Coaster
Conjecture (Conjecture~1.5 in the December 2016 preprint~\cite{LMroller})
prescribes both $q$ and $n=|V(G)|$. For every $q\ge2$ and $2q\le n\le3q-2$,
it asserts that all strict orderings of the coefficients indexed by
\[
 J(q,n)=\{\lceil q/2\rceil,\ldots,\lceil(n-1)/3\rceil\}
\]
occur among well-covered graphs of those parameters. The classical
Cutler--Pebody theorem~\cite{CProller} allows the order to vary.

\subsection{Two extensions of one face}
Consider a pure simplicial complex on $n$ vertices whose facets have
cardinality $q$. Write $s_i$ for the number of its faces with $i$ vertices.
For a $k$-face $S$, let $d_S$ be the number of vertices that can be added
individually to $S$, and let $e_S$ count the unordered pairs that can be
added together. Put $p=q-k$, $u=n-k$ and assume $0\le k\le q-2$.

\begin{lemma}\label{lem:face-slack}
The three consecutive face counts satisfy
\begin{align}
 &(p+u-1)(k+1)s_{k+1}-pu s_k-(k+1)(k+2)s_{k+2}\notag\\
 &\qquad=\sum_{|S|=k}\left[(d_S-p)(u-d_S)
          +2\left(\binom{d_S}{2}-e_S\right)\right]\ge0.
 \label{eq:face-slack}
\end{align}
Equality holds precisely when each $d_S$ is an endpoint of $[p,u]$ and
every pair of individually available vertices extends $S$ jointly.
\end{lemma}
\begin{proof}
Purity puts each $S$ in a facet, hence $p\le d_S\le u$.
Counting flags gives $\sum d_S=(k+1)s_{k+1}$ and
$\sum e_S=\binom{k+2}{2}s_{k+2}$. The identity follows upon expanding the
summand. Both terms on the right are nonnegative.
\end{proof}

\begin{theorem}\label{thm:fixedroller}
The prescribed-order conjecture is false. For each $k\ge5$, the admissible
pair $(q,n)=(2k,4k)$ excludes both strict orderings of the three ranks
$k,k+1,k+2$ in which the middle rank is smallest.
\end{theorem}
\begin{proof}
The independence complex of a well-covered graph is pure. For $q=2k$,
inequality~\eqref{eq:face-slack} gives, whenever $s_k\ge s_{k+1}$,
\[
 s_{k+2}\le \left(1-\frac{4k+3-n}{(k+1)(k+2)}\right)s_{k+1}.
\]
At $n=4k$ this is strictly smaller than $s_{k+1}$. For $k\ge5$,
$\lceil(4k-1)/3\rceil\ge k+2$, so all three ranks belong to $J(2k,4k)$.
Any permutation placing rank $k+1$ before the other two contradicts this
inequality. The parameter classes themselves are nonempty, as witnessed
by the disjoint union of $2k$ edges.
\end{proof}

The first member has $(q,n)=(10,20)$ and yields
$42s_7\le114s_6-75s_5$. A second, smaller witness is $(q,n)=(6,14)$,
where $20s_5\le52s_4-33s_3$. For this latter pair an older theorem already
gives a stronger one-step inequality. The Fisher--Ryan clique inequalities
applied to the complement~\cite{FisherRyan} give
$s_{j+1}/s_j\le n(q-j)/(q(j+1))$, by comparing successive normalized
roots and bounding each by $s_1/q$. Thus $s_5/s_4\le14/15$ at $(6,14)$.
In contrast, for the family $(2k,4k)$ that scalar upper bound at $j=k+1$
is $2(k-1)/(k+2)>1$.

\paragraph{Comparison with shadow inequalities.}
Take a balanced complex on twelve color classes with two vertices per
class, allowing every choice of at most one vertex of each color. In its
subcomplex $X_0$ retain both vertices in nine classes and one in the other
three. The face polynomial of $X_0$ is $(1+2x)^9(1+x)^3$, with counts
$24\,192,34\,848,36\,288$ at ranks $6,7,8$. Adjoin every ambient face of
size at most five and any $10\,657$ previously absent six-faces. There
are $2^6\binom{12}{6}=59\,136$ ambient six-faces, so these choices exist.
The resulting balanced complex $X$ has rank twelve on 24 vertices, and
\[
                 s_6=34\,849>34\,848=s_7< s_8=36\,288.
\]
As an actual balanced complex, $X$ satisfies all the
Frankl--F\"uredi--Kalai colored shadow inequalities and the ordinary
Kruskal--Katona inequalities. It also satisfies every first-order
purity inequality $(12-j)s_j\le(j+1)s_{j+1}$ of Michael and Traves
\cite{MTroller}. Indeed, $s_j/\binom{12}{j}=2^j$ through $j=5$, and
$32<34\,849/924<34\,848/792=44$. From rank seven onwards the counts
are those of the pure complex $X_0$, to which the first-order inequalities
apply by counting single extensions. Yet~\eqref{eq:face-slack} at
$(q,n,k)=(12,24,6)$ requires $56s_8\le161s_7-108s_6$.
Thus the three-rank restriction adds information even after these numerical
conditions are imposed together.

\subsection{Unimodality through independence number 35}
Call a well-covered graph \emph{very well-covered} if it has no isolated
vertices and $n=2q$. Levit and Mandrescu proved unimodality in this class
for $q\le9$~\cite{LMvery}. The following bound uses the compatibility of
extensions in the link of each independent set.
\begin{proposition}\label{prop:manyrank}
For a very well-covered graph of independence number $q$, fix
$0\le k<q$, put $p=q-k$, and take $1\le r\le p$. Then
\[
 \binom{k+r}{r}s_{k+r}\le
 \binom{p}{r}\left[s_k+
 \frac{2^r-1}{p}\big((k+1)s_{k+1}-p s_k\big)\right].
\]
\end{proposition}
\begin{proof}
Fix a perfect matching, whose existence follows from Favaron's
characterization~\cite{Favaron}. An independent $k$-set $S$ meets exactly
$k$ matched pairs. It extends to a maximum independent set, which meets
every matched pair once. Each of the $p$ unused pairs therefore offers at
least one vertex individually compatible with $S$. If $t$ pairs offer
both choices, then $d_S=p+t$. The number of compatible $r$-extensions
of $S$ is at most
\[
      h_r(t)=[x^r](1+2x)^t(1+x)^{p-t}.
\]
This counts all choices from distinct pairs and allows any additional
edges to reduce the count. For $0\le t\le p-2$, the second forward
difference of $h_r(t)$ equals
$[x^r]x^2(1+2x)^t(1+x)^{p-t-2}\ge0$. Its values thus lie below the
chord between $h_r(0)=\binom pr$ and $h_r(p)=2^r\binom pr$.
This chord bound also holds for $p=1$. Now sum over $S$, use
$\sum t=(k+1)s_{k+1}-p s_k$, and count the
$\binom{k+r}{r}$ starting $k$-sets in every $(k+r)$-set.
\end{proof}

\begin{theorem}\label{thm:very35}
Every very well-covered graph with $\alpha(G)\le35$ has a unimodal
independence polynomial.
\end{theorem}
\begin{proof}
Take $r=2$ in Proposition~\ref{prop:manyrank} and multiply by two to obtain
\[
 (k+1)(k+2)s_{k+2}
 \le(p-1)\{3(k+1)s_{k+1}-2p s_k\}.
\]
A weak descent $s_k\ge s_{k+1}$ therefore propagates one rank whenever
\[
 Q_q(k)=(k+1)(k+2)-(p-1)\{3(k+1)-2p\}
       =6k^2+(11-7q)k+2q^2-5q+5\ge0.
\]
The quadratic has discriminant $q^2-34q+1$, negative for $1\le q\le33$.
For $q=34$ the minimum over the reals is $-1/24$, and integral arguments
give integral values. For $q=35$ one has $Q_{35}(k)=6(k-19)(k-20)\ge0$
for every integer $k$. Thus every weak descent propagates through the
end, proving unimodality.
\end{proof}

\begin{corollary}\label{cor:very33}
Every very well-covered graph of independence number at most $33$ has a
log-concave independence sequence.
\end{corollary}
\begin{proof}
Set $x=s_{k+1}/s_k$. The two-extension bound in the preceding proof gives
$s_{k+1}^2\ge s_k s_{k+2}$ whenever
\[
 (k+1)(k+2)x^2-3(p-1)(k+1)x+2p(p-1)\ge0.
\]
For $p\ge2$, the discriminant is
$(p-1)(k+1)((p-9)(k-7)-72)$.
The two integer factors $p-9$ and $k-7$ have sum $q-16$, whose absolute
value is at most $17$ for $2\le q\le33$. Their product is at most
$17^2/4$, hence at most $72$ by integrality. The quadratic is therefore
nonnegative. The case $q=1$ has no triples of coefficients to compare.
\end{proof}

\subsection{Exact classification of the even-rank windows}

Put $h=\lceil\alpha/2\rceil$ and $u=\lceil(n-1)/3\rceil$. The complete
classification below concerns the assertion that \emph{every} strict order on
$\{h,\ldots,u\}$ is realized. It need not determine every feasible permutation.

\paragraph{Even independence number.}
For every admissible $n$ and every even $\alpha$, all permutations are
realizable if and only if $u-h+1\le2$.
Indeed, suppose $u\ge h+2$. If $n\le2\alpha+4$, the three-rank
secant bound at $k=h=\alpha/2$ excludes a strict valley. In detail,
write $A=s_k,B=s_{k+1},C=s_{k+2}$ and let $d$ be the number of
one-vertex extensions of a uniformly chosen $k$-face. Then
$k\le d\le n-k$ and $\mathbb E d=(k+1)B/A$.
Tur\'an's bound on the rank-$k$ link gives
$(k+1)(k+2)C/A\le (1-1/k)\mathbb E d^2$. The secant
$d^2\le nd-k(n-k)$ shows
\[
 (k+2)C/B\le (1-1/k)\left(n-\frac{k(n-k)}{\mathbb E d}\right).
\]
If $B\le A$, replace $\mathbb E d$ by its upper bound $k+1$.
Then $C<B$ follows from $(k-1)n<4k^2+2k$, which holds for
$n\le4k+4$. For $n\ge2\alpha+5$ put $t=u-1=\lfloor(n-2)/3\rfloor$.
We have
\[
t\ge(n-4)/3\ge\alpha(n-1)/(n+\alpha),
\qquad
\frac{s_{t+1}}{s_t}\le\frac{n(\alpha-t)}{\alpha(t+1)}\le1.
\]
The first line uses $n(n-2\alpha-4)\ge\alpha$; the ratio bound is the
classical graph Maclaurin inequality applied to the complement.

For a one-rank window existence follows from a disjoint union of $\alpha$
nonempty cliques of total order $n$. For a two-rank window write
$r=n-2\alpha$ and let $E_j$ denote the edgeless graph on $j$ vertices.
The join $E_\alpha\vee (K_{r+1}\sqcup E_{\alpha-1})$ is well-covered and has
polynomial $2(1+x)^\alpha+rx(1+x)^{\alpha-1}-1$, giving $s_h>s_{h+1}$.
The admissible pairs with even $\alpha$ and $u=h+1$ are
$(4,8),(4,9),(4,10),(6,12),(6,13),(8,16)$. In the same order,
$(\alpha-r)K_2\sqcup rK_3$ gives
$(s_h,s_{h+1})=(24,32),(30,44),(37,60),(160,240),(200,320),(1120,1792)$.

\paragraph{Exact three-rank spectra.}
For $(\alpha,n)=(6,14)$ the graph Maclaurin ratio
$s_5/s_4\le14/15$ restricts the six strict orders to three. All three
are attained even among connected graphs. Take five blocks $K_2$ and one
block $K_4$. In each block mark one vertex, and place a graph $F$ on the six
marks. Every maximal independent set meets every block because each contains
an unmarked vertex with no external neighbors. If $F$ is a star centered in
the block $K_4$, the polynomial is
$(1+3x)(1+2x)^5+x(1+x)^5$ and the triple $(s_3,s_4,s_5)$ is $(210,330,277)$.
If $F=K_6$ the polynomial is
$(1+3x)(1+x)^5+x(1+x)^5+5x(1+3x)(1+x)^4$, giving $(140,155,86)$.
Finally, for $H=5K_1\sqcup K_2$, the join $H\vee H$ gives
$2(1+x)^5(1+2x)-1$ and $(60,50,22)$.

For $(\alpha,n)=(6,16)$ exactly four orders are realized. The
two strict valleys are forbidden by the same secant at $k=3$. Witnesses
for the other four are as follows:
\[
\begin{array}{lrrr}
G &s_3&s_4&s_5\\
2K_2\sqcup4K_3 &372&729&756\\
5K_2\sqcup K_6 &320&560&512\\
\text{marked clique on the six blocks of }5K_2\sqcup K_6 &200&235&136\\
H\vee H,\quad H=4K_1\sqcup2K_2 &88&82&40
\end{array}
\]
The last two polynomials are
$(1+5x)(1+x)^5+x(1+x)^5+5x(1+5x)(1+x)^4$ and
$2(1+x)^4(1+2x)^2-1$. Each witness is well-covered and has no isolated vertices.



\section{A chordal bipartite counterexample on 1,408 vertices}
\label{sec:bipartite-valley}
The unrestricted very well-covered graph conjecture already has Baker's
4,074-vertex counterexample~\cite{Baker}. Its core contains cliques on twelve vertices.
The construction below gives an infinite family with no induced cycle of
length greater than four, and one explicit member of order 1,408.
Bhattacharyya and Kahn previously obtained nonunimodality for bipartite
graphs using a complete bipartite core and a matching on one side
\cite{BK}. Their two sides are maximal
independent sets of different sizes. The balanced blocks below make
every maximal independent set have the same size while preserving the
two separated coefficient peaks.

For positive integers $a,b$, take $a$ disjoint edges $x_i y_i$ and $b$
disjoint copies of $K_{4,4}$ with partite classes $U_j,V_j$. Add every
edge $x_i u$ with $u\in\bigcup_j U_j$, and call the resulting graph
$H(a,b)$. It is connected, with bipartition
$\{x_i\}_i\cup\bigcup_j V_j$ and $\{y_i\}_i\cup\bigcup_j U_j$.
Every maximal independent set contains exactly one of $x_i,y_i$ for each
$i$, since $y_i$ has only the neighbor $x_i$. It contains exactly one
whole side of each block. Indeed, if it avoids $U_j$ then all of $V_j$
can be added, while if it meets $U_j$ it avoids all $x_i$ and all of
$V_j$, so that all of $U_j$ can be added. Thus $H(a,b)$ is very
well-covered with independence number $N=a+4b$.

Any induced cycle avoiding all $x_i$ lies in a single block, and hence
has length four. For an induced cycle containing some $x_i$, all its
vertices in $\bigcup U_j$ must be the two neighbors of $x_i$ on the cycle,
since any further one would give a chord. No $y_i$ lies on a cycle,
so this cycle also has length four by bipartiteness. Thus $H(a,b)$ is
chordal bipartite.

\begin{theorem}\label{thm:explicit-bipartite-valley}
The independence polynomial of $H(204,125)$ has
$s_{373}>s_{374}<s_{375}$. In particular, unimodality can fail on
connected chordal bipartite very well-covered graphs.
\end{theorem}
\begin{proof}
An independent set either avoids all $U_j$ or avoids all $x_i$.
The generating polynomials of these two classes and their intersection give
\[
 I(H(a,b);t)=(1+t)^{4b}(1+2t)^a
    +(1+t)^a\bigl(2(1+t)^4-1\bigr)^b-(1+t)^N.
\]
For $(a,b)=(204,125)$, exact integer expansion gives
\[
 10^{242}<s_{373}-s_{374}<2\cdot10^{242},\qquad
 10^{241}<s_{375}-s_{374}<2\cdot10^{241}.
\]
The entire arithmetic check is reproduced below in Python 3. The
convolution uses $2(1+t)^4-1=1+8t+12t^2+8t^3+2t^4$.
\begin{lstlisting}
from math import comb
A, B = 204, 125
v = [comb(A, i) for i in range(A+1)]
for _ in range(B):
    w = [0] * (len(v)+4)
    for i, z in enumerate(v):
        for j, q in enumerate((1, 8, 12, 8, 2)):
            w[i+j] += q*z
    v = w
def s(k):
    first = sum(comb(A,i)*2**i*comb(4*B,k-i)
                for i in range(max(0,k-4*B), min(k,A)+1))
    return first + v[k] - comb(A+4*B,k)
assert 10**242 < s(373)-s(374) < 2*10**242
assert 10**241 < s(375)-s(374) < 2*10**241
\end{lstlisting}
\end{proof}

\begin{theorem}\label{thm:bipartite-infinite}
Put $\lambda=\log(31/16)/\log(3/2)$ and $a=\lfloor\lambda b\rfloor$.
For all sufficiently large $b$, the independence polynomial of $H(a,b)$
fails to be unimodal.
\end{theorem}
\begin{proof}
Denote the first two positive summands of the displayed counting formula
by $A_b,B_b$, and the subtracted polynomial by $C_b$. At $t=1$ their
masses satisfy
\[
 2/3<A_b(1)/B_b(1)\le1,\qquad C_b(1)/B_b(1)=(16/31)^b.
\]
Normalize each positive summand by its mass. Its coefficients give the
distribution of a sum of independent integer-valued random variables,
one per factor. Logarithmic differentiation at $1$ gives the means
\[
 \mu_A=N/2+a/6,\qquad \mu_B=N/2+2b/31 .
\]
Their difference grows linearly in $b$, since $31/16>3/2$ implies
$\lambda>1$. Both variances are $O(b)$, as every factor has degree
at most four and there are $O(b)$ factors. Chebyshev's inequality gives
each distribution mass at least $1/2$ within $K\sqrt b$ of its mean
for a fixed $K$. Each of these intervals therefore contains an index
with normalized coefficient at least a constant times $b^{-1/2}$.
At the nearest integer to the midpoint of the means, each distribution
has probability $O(b^{-1})$, again by Chebyshev. After multiplying by
the respective masses, the two intervals contain coefficients of
$I(H(a,b);t)$ of order at least $B_b(1)b^{-1/2}$, since each coefficient
of $C_b$ is at most $C_b(1)$. The coefficient at the intervening midpoint
is $O(B_b(1)/b)$. This strict valley across three ordered indices proves
the assertion.
\end{proof}

\begin{theorem}\label{thm:bipartite-many-peaks}
For each integer $m\ge1$ a connected chordal bipartite very well-covered graph
has an independence sequence with $m+1$ peak regions separated by valleys.
Precisely, there exist indices
$h_0<\ell_1<h_1<\cdots<\ell_m<h_m$ such that
$s_{\ell_j}<\min\{s_{h_{j-1}},s_{h_j}\}$ for every $j$.
\end{theorem}
\begin{proof}
Use $m$ copies of $H(\lfloor\lambda b\rfloor,b)$ with
$b$ sufficiently large for the fixed $m$. Choose one $x_i$, denoted $x_*$,
in the first copy and join it to all four members of $U_1$ in every other
copy. These edges respect the bipartition and keep the graph connected.
The earlier maximality argument remains valid: the vertices in each $U_j$
still have identical neighborhoods, and the $V_j$ have no external neighbors.
Deleting $x_*$ separates the other copies from the first and from each other.
An induced cycle using a new edge therefore uses $x_*$ and lies otherwise
in one other copy. Its two neighbors at $x_*$ are in the same $U_1$.
These two vertices have identical neighborhoods, so a cycle containing
both at distance two has a chord unless its length is four.

Use $A_b,B_b$ from the preceding proof. The resulting independence
polynomial is bounded coefficientwise above by $(A_b+B_b)^m$.
For each $i=0,\ldots,m$, prescribe the class represented by $A_b$ in
exactly $i$ fixed copies and that represented by $B_b$ in the others.
Also forbid $x_*$. All the resulting sets are independent in the
connected graph. If the first copy is in the $A_b$-class, this last
restriction replaces one $(1+2t)$ factor by $(1+t)$; in the $B_b$-class
it is automatic. Each prescribed class therefore has mass comparable to
$B_b(1)^m$, variance $O_m(b)$ and mean within a fixed constant of
$M_i=i\mu_A+(m-i)\mu_B$.
Consecutive $M_i$ differ by a positive constant times $b$ up to $O(1)$.
Chebyshev's inequality gives an index $h_i$ within $O_m(\sqrt b)$ of
$M_i$ where the graph coefficient is at least a constant times
$B_b(1)^m/\sqrt b$. At a nearest integer $\ell_j$ to the midpoint of
$M_{j-1}$ and $M_j$, every one of the $m+1$ terms in the binomial expansion
of the upper bound contributes $O_m(B_b(1)^m/b)$, by the same variance
estimate. Taking $b$ large proves all $m$ strict inequalities.
\end{proof}
Bhattacharyya and Kahn also proved the existence of arbitrarily many
peaks for bipartite graphs. The preceding construction simultaneously
enforces equal size of every maximal independent set and gives an
explicit induced-cycle restriction.

\section{Divisibility in reciprocal Lucas recurrences}

Z.-W. Sun's Conjecture~4.4 in \emph{Some new problems in additive
combinatorics}~\cite{Sun2019} (posed September 27, 2014, published 2019) concerns
the integer sequences
\[
 U_0=0,\quad U_1=1,\qquad V_0=2,\quad V_1=a,\qquad
 X_{j+2}=aX_{j+1}+X_j.
\]
For each integer $a\not\equiv3\pmod6$ and each $m\ge1$, it asks
for infinitely many positive $n$ satisfying $m+n\mid U_m+U_n$,
and separately for infinitely many positive $n$ satisfying
$m+n\mid V_m+V_n$. The two sets of witnesses may differ. For the related problem of divisibility of a single Lucas term by its own index, Smyth classified the admissible indices~\cite{Smyth}. The shifted sums here require a prescribed sign at the complementary index $N-m$.

\paragraph{Classification.}
Within the stipulated parameter set, both assertions hold for every
$m\ge1$ precisely when $a\ne0$. At $a=0$ the first assertion holds
precisely for even $m$, and the second precisely for odd $m$.
For every nonzero parameter, the proof produces geometric progressions
of moduli $m+n$ for each assertion.

Indeed, at $a=0$ the two sequences alternate between $0,1$ and $2,0$,
respectively. A vanishing sum gives infinitely many witnesses, whereas
a nonzero sum is bounded by $4$.

\subsection{Two matrix congruences}

Put $A=\left(\begin{smallmatrix}a&1\\1&0\end{smallmatrix}\right)$.
Since $\det A=-1$, every integral power, positive or negative, has
integer entries. For each $j$, its upper right entry is $U_j$ and its
trace is $V_j$. Inversion of a two-by-two matrix gives
\[
 U_{-m}=(-1)^{m+1}U_m,\qquad V_{-m}=(-1)^m V_m .
\]
Given the sequence and $m$, choose $s=\pm1$ so that
$X_m+sX_{-m}=0$ as an integer. Then $N>m$ and
$A^N\equiv sI\pmod N$ suffice, since $n=N-m$.
One can also verify this congruence on an individual entry or trace
at one factor of $N$.

We repeatedly use two elementary lifting statements. If a matrix
$B\equiv sI\pmod p$ where $s=\pm1$, then
$B^{p^k}\equiv s^{p^k}I\pmod{p^{k+1}}$. More generally, if $B$ is
scalar modulo $p$, then $B^{p^k}$ is scalar modulo $p^{k+1}$.
Both statements follow by writing $B=cI+pH$ and iterating the binomial
expansion of the $p$th power. The terms linear in the error gain a
factor $p$, and those of degree at least two already gain a full
additional power.

For $s=+1$ and $a$ even use $N=2^k$, $k\ge1$.
The identity $A^2=aA+I$ gives $A^2\equiv I\pmod2$.
If $a$ is odd and admissible then $3\nmid a$. Use $N=3\cdot2^k$,
$k\ge3$. The identity $A^3=(a^2+1)A+aI$ gives $A^3\equiv I\pmod2$.
Modulo $3$, $A^4=(a^3+2a)A+(a^2+1)I\equiv-I$, so $A^8\equiv I$.
Lifting at $2$ and reducing at $3$ give $A^N\equiv I\pmod N$.

For $s=-1$, first let $p$ be any odd prime divisor of $a^2+4$.
Modulo $p$, put $c=a/2$. Then $A=cI+J$ with $J^2=0$ and $c^2=-1$.
It follows that $A^{2p}=-I$ in characteristic $p$.
Lifting at $p$ gives $A^{2p^k}\equiv-I\pmod{p^k}$ for every $k\ge1$.
If $a=2b$ with $|b|>1$, such an odd prime exists: when $b$ is even,
$b^2+1$ is odd and greater than one, and when $b$ is odd, $b^2+1$ is
$2\pmod8$ and greater than two. In this case $N=2p^k$ also gives
$A^N\equiv I\equiv-I\pmod2$, completing the construction.

Let $a$ be odd and admissible. If $3\nmid m$, choose $p\mid a^2+4$
with $p\equiv2\pmod3$. This is possible since $a^2+4$ is both odd
and $2\pmod3$. Choose arbitrarily large $k$ of the parity for which
$p^k\equiv m\pmod3$. Put $N=2p^k$. Modulo $p^k$ the sign is negative.
Modulo $2$ we have $N-m\equiv m\pmod3$ and $A^3=I$, so the two
terms have the same parity. This proves the divisibility modulo $N$.
If $3\mid m$, take $N=4\cdot3^k$, $k\ge1$. From $A^4\equiv-I\pmod3$
we obtain $A^N\equiv-I\pmod{3^k}$. Modulo $4$, the identity
$A^3\equiv2A+aI$ gives $A^6\equiv I$. Thus $A^N\equiv I\pmod4$.
Since $3\mid m$ and $A^3\equiv I\pmod2$, both the upper right entry
and the trace of $A^{-m}$ are even. Their positive and negative values
therefore agree modulo $4$. This supplies the required negative sign
on each tested quantity.

\subsection{The Pell family}

It suffices to handle $a=2$. Indeed,
$U_j(-2)=(-1)^{j-1}U_j(2)$ and $V_j(-2)=(-1)^jV_j(2)$ for $j\ge1$,
and all the displayed moduli are even. The positive-sign construction
applies as above. For $V$ and the negative sign, take $N=4\cdot3^k$.
Here $A^4\equiv-I\pmod3$, while $A^2=I+2A$ gives $A^4\equiv I\pmod4$.
All $V_j$ are even, so once more the two signs agree modulo $4$ on the
trace.

It remains to treat $U$ with $m$ odd. For odd $r>m$, set $n=2r-m$.
The addition identity
\[
 U_m+U_{2r-m}=U_r V_{r-m}
\]
follows from the characteristic roots $\alpha,\beta$ via
$\alpha\beta=-1$ and the fact that $r-m$ is even. For integers $d\ge1$
and odd $h\ge1$, $V_d\mid V_{dh}$ because $x+y$ divides $x^h+y^h$,
and the symmetric quotient evaluated at $x=\alpha^d,y=\beta^d$
is an integer. If a prime $b\mid U_h$, then $A^h$ is scalar modulo
$b$, by $A^h=U_h A+U_{h-1}I$. Hence $b^k\mid U_{h b^k}$.

If $m\equiv1\pmod4$, take $r=3\cdot5^k$. As $U_3=5$, we have
$5^k\mid U_r$. Since $r-m\equiv2\pmod4$, $V_2=6$ divides
$V_{r-m}$, so $2r\mid U_r V_{r-m}$.

If $m\equiv3\pmod4$, let $r=57b^k$, with $57=3\cdot19$.
For $5\mid m$, choose $b=5$ and arbitrary $k\ge1$.
Otherwise choose $b=37$ and take all sufficiently large $k$
in a residue class modulo four for which $57\cdot37^k\equiv m\pmod5$.
Such a class exists because $37\equiv2\pmod5$ generates the
multiplicative group modulo five. In either case $r\equiv1\pmod4$
and $r-m\equiv10\pmod{20}$. Consequently $V_{10}=6726=2\cdot3\cdot19\cdot59$
divides $V_{r-m}$. The two candidate primes divide $U_{57}$: indeed
$U_3=5$, $U_{19}=6625109=37\cdot179057$, and the identity
$A^h=U_hA+U_{h-1}I$ shows that $U_{dh}\equiv0\pmod{U_h}$.
The lifting observation therefore gives $b^k\mid U_r$.
Since $b\notin\{2,3,19\}$, we get $2r\mid U_r V_{r-m}$.
The available exponents are unbounded in every case.

\paragraph{Fibonacci and Lucas witnesses.}
For $a=1$ we have $U_j=F_j$ and $V_j=L_j$. Thus, for each $m\ge1$,
the sets
\[
 \{n>m:m+n\mid F_m+F_n\},\qquad
 \{n>m:m+n\mid L_m+L_n\}
\]
are both infinite. Their minima, which define Sun's sequences
A247937 and A247940, respectively, therefore exist at every index
\cite{SunOEISF,SunOEISL}. The Fibonacci construction for $m=1$,
for example, supplies every modulus $N=2\cdot5^{2j}$ with $j\ge1$
and $n=N-1$.

\paragraph{Bounds for the number and first occurrence of witnesses.}
For each fixed nonzero admissible $a$ there is a constant $C_a$ such
that for either sequence and every $m\ge1$ there is a witness
$m<n\le C_a m$.
Fix, once for each parameter in the relevant case, an eligible prime
$p\mid a^2+4$. Each branch above supplies a progression
$N=B R^j$ ($j\ge0$) of integer moduli with $R>1$.
The choice of branch depends only on finitely many residue classes of
$m$ for each $a$. In a Pell branch one can require $N>2m$, as in its
proof. In every branch the first modulus exceeding $2m$ is at most
$\max\{B,2Rm\}$: if it is not $B$, its predecessor is at most $2m$.
Taking the maximum of $B$ and $2R$ over the finitely many progressions
proves the asserted linear bound. Moreover, counting the same moduli
in $(2m,X+m]$ shows that for constants $c_a>0$ and $K_a$ independent
of $m$ the number of witnesses $m<n\le X$ is at least
$c_a\log(X/m)-K_a$ for $X\ge m$.
Indeed, write $B_* =\max\{B,2R\}$ across the progressions,
and $R_* =\max R$. Starting at a modulus at most $B_* m$,
the first $1+\lfloor\log((X+m)/(B_*m))/\log R_*\rfloor$ terms
are all counted whenever $(X+m)/(B_*m)\ge1$. Taking
$c_a=1/\log R_*$ and $K_a=\log B_*/\log R_*$ handles smaller $X$
as well.

When $a=1$, the positive-sign progression starts at $24$ with ratio $2$.
The negative-sign progression for $3\mid m$ starts at $12$ with
ratio $3$. For $3\nmid m$ it starts at $10$ or $50$ with ratio $25$.
Thus $B\le50$ and $R\le25$ in every branch. Since $m\ge1$,
$\max\{50,50m\}=50m$ even at the initial indices. Each of the two
classical OEIS sequences consequently has its $m$th entry at most $49m$.

\section{What the results establish}
The source explicitly records the $d=2$ solution of Open Problem~1. Theorem~\ref{thm:sampling} treats every remaining dimension and is uniform over the original normalized class. For Open Problems~2 and 3 the source had treated $d=2,3,5,7$. The additional dimensions evaluated here are $4,6$ and $8$ through $32$. The exact maximizing index at each individual $d\ge33$ remains to be classified. Equations~\eqref{eq:bb-tail} and~\eqref{eq:bb-limit} supply finite candidates and a limiting law in this range.

For the fixed-order graph conjecture Fisher--Ryan already implies a counterexample at $(6,14)$. The simultaneous three-rank constraint, the even-rank classification and the bounds in Theorem~\ref{thm:very35} and Corollary~\ref{cor:very33} are the additional conclusions proved here. The unrestricted very well-covered unimodality counterexample and the two-term bipartite construction have prior precedents~\cite{Baker,BK}. In the arithmetic problem Sun's printed parameter range includes zero, so the stated universal conjecture is false. The nonzero-parameter theorem gives full proofs of the Fibonacci and Lucas cases, with a quantitative witness bound.

The original statements and their applicable quantifiers are identifiable in the cited texts. Our comparison with accessible literature and the supplied corpus did not locate an earlier full solution to the coefficient summability problem or the same dimension classification. Historical priority for these conclusions and for the method remains provisional. Each use of Fourier representation, separated sampling, Sturm comparison, graph clique bounds, or matrix lifting is compatible with earlier work on those tools.

\appendix
\section{Exact certificate for the maximizing indices}
This Python~3 program uses integer and rational arithmetic throughout. The finite checks, together with \eqref{eq:bb-tail}, prove the strict global comparisons in Section~\ref{sec:bb-ratios}. The output gives the seven integer parts tabulated there.
\begin{lstlisting}
from fractions import Fraction
from math import prod


def enclosure(order, x, kind):
    x = Fraction(x)
    term = x**order / prod(range(1, 2*order + 2, 2))
    total = Fraction(0)
    for k in range(147):
        power = order + 2*k
        value = term
        if kind == 'derivative':
            value = term * power / x
        if kind == 'integral':
            value = term * x / (power + 1)
        if k == 146:
            return min(total, total + value), max(total, total + value)
        total += value
        term *= -x*x / (2*(k + 1)*(2*order + 2*k + 3))


def sign(order, x):
    lower, upper = enclosure(order, x, 'function')
    assert lower * upper > 0
    return 1 if lower > 0 else -1


def peak_interval(order, left):
    a, b = Fraction(left), Fraction(left + 1)
    for _ in range(18):
        mid = (a + b)/2
        if sign(order, mid) > 0:
            a = mid
        else:
            b = mid
    mid, error = (a + b)/2, (b - a)/2
    v, w = enclosure(order, mid, 'derivative')
    y, z = enclosure(order, mid, 'integral')
    assert y > error and -w > error/3
    return (-w - error/3)/(z + error), (-v + error/3)/(y - error)


mass = [Fraction(157, 100), Fraction(1)]
for order in range(2, 31):
    mass.append(mass[order-2] * Fraction(order-1, order))
gaps = []
for order in range(31):
    intervals = []
    start = max(1, order)
    last = sign(order, start)
    assert last == 1
    for left in range(start, 70):
        right = sign(order, left+1)
        if last == 1 and right == -1:
            intervals.append(peak_interval(order, left))
        last = right
    winner = 0 if order <= 26 else 1
    rivals = [1/(70*mass[order])]
    rivals += [v[1] for i,v in enumerate(intervals) if i != winner]
    gap = intervals[winner][0] - max(rivals)
    assert gap > Fraction(1, 12000)
    gaps.append(gap)

ranges = [(0,5), (6,10), (11,15), (16,20), (21,25), (26,26), (27,30)]
certificate = [int(10**6 * min(gaps[a:b+1])) for a,b in ranges]
assert certificate == [45756, 21405, 10414, 4479, 918, 380, 101]
print(certificate)
\end{lstlisting}

\begin{thebibliography}{99}

\bibitem{BB14N} E. Berdysheva, Question on spherical Bessel functions, \emph{OP-SF Net} 21.3 (May 2014), Topic 4. \href{https://math.nist.gov/~DLozier/OPSFnet/OPSF-Net_21-3.pdf}{NIST archive}.
\bibitem{BB14} E. E. Berdysheva and H. Berens, On a discrete Tur\'an problem for $\ell_1$ radial functions, in \emph{New Perspectives on Approximation and Sampling Theory}, Birkh\"auser, 2014, 447--468. \href{https://doi.org/10.1007/978-3-319-08801-3_19}{Chapter}.
\bibitem{BB16} E. E. Berdysheva and H. Berens, On the Tur\'an problem for positive definite $\ell_1$-radial functions (in Russian), Stechkin conference proceedings, 2016, 65--72. Open Problems 1--3 appear on pp.~70--71. \href{https://ir.duan.edu.ua/server/api/core/bitstreams/fdaadc5f-8619-438f-a679-98a574d5aaca/content}{Proceedings}.
\bibitem{deboor} C. de Boor, Divided differences, \emph{Surveys in Approximation Theory} \textbf{1} (2005), 46--69. \href{https://arxiv.org/abs/math/0502036}{Preprint}.
\bibitem{BBDLMF} NIST, \emph{Digital Library of Mathematical Functions}, Chapter 10, Bessel functions. \url{https://dlmf.nist.gov/10}.
\bibitem{LMroller} V. E. Levit and E. Mandrescu, The Roller-Coaster Conjecture revisited, \emph{Graphs Combin.} \textbf{33} (2017), 1499--1508. Conjecture numbering follows the \href{https://arxiv.org/abs/1612.03736}{December 2016 preprint}.
\bibitem{CProller} J. Cutler and L. Pebody, Maximal-clique partitions and the Roller Coaster Conjecture, \emph{J. Combin. Theory Ser. A} \textbf{145} (2017), 25--35. \href{https://doi.org/10.1016/j.jcta.2016.06.019}{Article}.
\bibitem{FisherRyan} D. C. Fisher and J. Ryan, Bounds on the number of complete subgraphs, \emph{Discrete Math.} \textbf{103} (1992), 313--320. \href{https://doi.org/10.1016/0012-365X(92)90323-8}{Article}.
\bibitem{MTroller} T. S. Michael and W. N. Traves, Independence sequences of well-covered graphs: non-unimodality and the Roller-Coaster conjecture, \emph{Graphs Combin.} \textbf{19} (2003), 403--411.
\bibitem{LMvery} V. E. Levit and E. Mandrescu, Very well-covered graphs and the unimodality conjecture, 2004. \url{https://arxiv.org/abs/math/0406623}.
\bibitem{Favaron} O. Favaron, Very well covered graphs, \emph{Discrete Math.} \textbf{42} (1982), 177--187. \href{https://doi.org/10.1016/0012-365X(82)90215-1}{Article}.
\bibitem{Baker} L. Baker, a very well-covered graph with a nonunimodal independence polynomial, public announcement (July 2026). \href{https://aiproofwatch.com/updates/a-very-well-covered-graph-with-a-nonunimodal-independence-polynomial-bc5e7b0b}{Exact check and link to announcement}.
\bibitem{BK} A. Bhattacharyya and J. Kahn, A bipartite graph with non-unimodal independent set sequence, \emph{Electron. J. Combin.} \textbf{20} (2013), P11. \href{https://doi.org/10.37236/3034}{Article}.
\bibitem{VMRC} QEEDDD, \emph{Verifiable Mathematical Reasoning Corpus}, entries 185, 199, 204, and 216. \href{https://github.com/QEEDDD/Verifiable-Mathematical-Reasoning-Corpus/tree/c2ac34ea4622778187173f41bb1996deee8081b7}{Repository snapshot, c2ac34e}.
\bibitem{Smyth} C. Smyth, The terms in Lucas sequences divisible by their indices, \emph{J. Integer Seq.} \textbf{13} (2010), Article 10.2.4. \href{https://cs.uwaterloo.ca/journals/JIS/VOL13/Smyth/smyth2.html}{Article}.

\bibitem{Sun2019}
Z.-W. Sun, Some new problems in additive combinatorics,
\emph{Nanjing University Journal of Mathematics. Biquarterly}
\textbf{36} (2019), no. 2, 134--155, Conjecture 4.4.
\url{https://arxiv.org/abs/1309.1679}.
\bibitem{SunOEISF}
Z.-W. Sun, Sequence A247937, \emph{The On-Line Encyclopedia of
Integer Sequences}, September 27, 2014. \url{https://oeis.org/A247937}.
\bibitem{SunOEISL}
Z.-W. Sun, Sequence A247940, \emph{The On-Line Encyclopedia of
Integer Sequences}, September 27, 2014. \url{https://oeis.org/A247940}.

\end{thebibliography}

\end{document}
